\documentclass[10pt,a4paper]{amsart}
\usepackage{amsmath,amssymb,amsthm,graphicx}
% For commutative diagrams we use Xy-pic
\usepackage[all]{xy}

% We use 2cell for 2-commutative diagrams.
\xyoption{2cell}
% Theorem environments.
%
\theoremstyle{plain}
\newtheorem{theorem}[subsection]{Theorem}
\newtheorem{proposition}[subsection]{Proposition}
\newtheorem{lemma}[subsection]{Lemma}

\theoremstyle{definition}
\newtheorem{definition}[subsection]{Definition}
\newtheorem{example}[subsection]{Example}
\newtheorem{exercise}[subsection]{Exercise}
\newtheorem{situation}[subsection]{Situation}

\theoremstyle{remark}
\newtheorem{remark}[subsection]{Remark}
\newtheorem{remarks}[subsection]{Remarks}

\numberwithin{equation}{subsection}

% Macros
%
\def\lim{\mathop{\mathrm{lim}}\nolimits}
\def\colim{\mathop{\mathrm{colim}}\nolimits}
\def\Spec{\mathop{\mathrm{Spec}}}
\def\Hom{\mathop{\mathrm{Hom}}\nolimits}
\def\Ext{\mathop{\mathrm{Ext}}\nolimits}
\def\SheafHom{\mathop{\mathcal{H}\!\mathit{om}}\nolimits}
\def\SheafExt{\mathop{\mathcal{E}\!\mathit{xt}}\nolimits}
\def\Sch{\mathit{Sch}}
\def\Mor{\mathop{\mathrm{Mor}}\nolimits}
\def\Ob{\mathop{\mathrm{Ob}}\nolimits}
\def\Sh{\mathop{\mathit{Sh}}\nolimits}
\def\NL{\mathop{N\!L}\nolimits}
\def\CH{\mathop{\mathrm{CH}}\nolimits}
\def\proetale{{pro\text{-}\acute{e}tale}}
\def\etale{{\acute{e}tale}}
\def\QCoh{\mathit{QCoh}}
\def\Ker{\mathop{\mathrm{Ker}}}
\def\Im{\mathop{\mathrm{Im}}}
\def\Coker{\mathop{\mathrm{Coker}}}
\def\Coim{\mathop{\mathrm{Coim}}}

% Boxtimes
%
\DeclareMathSymbol{\boxtimes}{\mathbin}{AMSa}{"02}

%
% Macros for moduli stacks/spaces
%
\def\QCohstack{\mathcal{QC}\!\mathit{oh}}
\def\Cohstack{\mathcal{C}\!\mathit{oh}}
\def\Spacesstack{\mathcal{S}\!\mathit{paces}}
\def\Quotfunctor{\mathrm{Quot}}
\def\Hilbfunctor{\mathrm{Hilb}}
\def\Curvesstack{\mathcal{C}\!\mathit{urves}}
\def\Polarizedstack{\mathcal{P}\!\mathit{olarized}}
\def\Complexesstack{\mathcal{C}\!\mathit{omplexes}}
% \Pic is the operator that assigns to X its picard group, usage \Pic(X)
% \Picardstack_{X/B} denotes the Picard stack of X over B
% \Picardfunctor_{X/B} denotes the Picard functor of X over B
\def\Pic{\mathop{\mathrm{Pic}}\nolimits}
\def\Picardstack{\mathcal{P}\!\mathit{ic}}
\def\Picardfunctor{\mathrm{Pic}}
\def\Deformationcategory{\mathcal{D}\!\mathit{ef}}
\newenvironment{reference}{}{}
\newenvironment{slogan}{}{}
\newenvironment{history}{}{}
\usepackage{cleveref}
\crefname{originalequation}{original equation}{original equations}
\Crefname{originalequation}{Original equation}{Original equations}
\makeatletter
\newlabel{schemes-equation-canonical-morphism}{{02NA}{source}{}{sourcepointer-1}{}}
\newlabel{schemes-equation-canonical-morphism@cref}{{[originalequation][0][]02NA}{[1][0][]source}{}{}{}}
\makeatother
\crefname{originallemma}{original lemma}{original lemmas}
\Crefname{originallemma}{Original lemma}{Original lemmas}
\makeatletter
\newlabel{spaces-cohomology-lemma-flat-base-change-cohomology}{{073K}{source}{}{sourcepointer-2}{}}
\newlabel{spaces-cohomology-lemma-flat-base-change-cohomology@cref}{{[originallemma][0][]073K}{[1][0][]source}{}{}{}}
\makeatother
\crefname{originaltheorem}{original theorem}{original theorems}
\Crefname{originaltheorem}{Original theorem}{Original theorems}
\makeatletter
\newlabel{spaces-cohomology-section-theorem-formal-functions}{{08AU}{source}{}{sourcepointer-3}{}}
\newlabel{spaces-cohomology-section-theorem-formal-functions@cref}{{[originaltheorem][0][]08AU}{[1][0][]source}{}{}{}}
\makeatother
\crefname{originallemma}{original lemma}{original lemmas}
\Crefname{originallemma}{Original lemma}{Original lemmas}
\makeatletter
\newlabel{spaces-cohomology-lemma-i-star-equivalence}{{08AM}{source}{}{sourcepointer-4}{}}
\newlabel{spaces-cohomology-lemma-i-star-equivalence@cref}{{[originallemma][0][]08AM}{[1][0][]source}{}{}{}}
\makeatother
\crefname{originallemma}{original lemma}{original lemmas}
\Crefname{originallemma}{Original lemma}{Original lemmas}
\makeatletter
\newlabel{coherent-lemma-graded-finiteness}{{0897}{source}{}{sourcepointer-5}{}}
\newlabel{coherent-lemma-graded-finiteness@cref}{{[originallemma][0][]0897}{[1][0][]source}{}{}{}}
\makeatother
\crefname{originallemma}{original lemma}{original lemmas}
\Crefname{originallemma}{Original lemma}{Original lemmas}
\makeatletter
\newlabel{spaces-cohomology-lemma-relative-affine-cohomology}{{0D2U}{source}{}{sourcepointer-6}{}}
\newlabel{spaces-cohomology-lemma-relative-affine-cohomology@cref}{{[originallemma][0][]0D2U}{[1][0][]source}{}{}{}}
\makeatother
\crefname{originaltheorem}{original theorem}{original theorems}
\Crefname{originaltheorem}{Original theorem}{Original theorems}
\makeatletter
\newlabel{spaces-cohomology-theorem-formal-functions}{{08AZ}{source}{}{sourcepointer-7}{}}
\newlabel{spaces-cohomology-theorem-formal-functions@cref}{{[originaltheorem][0][]08AZ}{[1][0][]source}{}{}{}}
\makeatother
\crefname{originallemma}{original lemma}{original lemmas}
\Crefname{originallemma}{Original lemma}{Original lemmas}
\makeatletter
\newlabel{spaces-cohomology-lemma-proper-over-affine-cohomology-finite}{{08AS}{source}{}{sourcepointer-8}{}}
\newlabel{spaces-cohomology-lemma-proper-over-affine-cohomology-finite@cref}{{[originallemma][0][]08AS}{[1][0][]source}{}{}{}}
\makeatother
\crefname{originallemma}{original lemma}{original lemmas}
\Crefname{originallemma}{Original lemma}{Original lemmas}
\makeatletter
\newlabel{algebra-lemma-NAK}{{00DV}{source}{}{sourcepointer-9}{}}
\newlabel{algebra-lemma-NAK@cref}{{[originallemma][0][]00DV}{[1][0][]source}{}{}{}}
\makeatother
\crefname{originallemma}{original lemma}{original lemmas}
\Crefname{originallemma}{Original lemma}{Original lemmas}
\makeatletter
\newlabel{spaces-cohomology-lemma-ML-cohomology-powers-ideal}{{08AY}{source}{}{sourcepointer-10}{}}
\newlabel{spaces-cohomology-lemma-ML-cohomology-powers-ideal@cref}{{[originallemma][0][]08AY}{[1][0][]source}{}{}{}}
\makeatother
\title{Uniform high-twist vanishing and section lifting from an ample scheme fibre of a proper algebraic space}
\author{The Stacks Project Authors}
\date{stacks-project:spaces-divisors.tex:local-snapshot-20261011; GNU FDL 1.2 or later}
\begin{document}
\maketitle
\noindent\textit{Editorial note.} The conclusion concerns stalks aty for alld>=d\_0. The fibreX\_y is expressly assumed to be a scheme; no global scheme assumption onX is introduced.\par\medskip
\noindent\textit{Allowed standard prerequisites (editorial).}\begin{itemize}
\item Explicit prerequisite boundary. Proper algebraic spaces over a Noetherian scheme, coherent sheaves, ampleness andSerre vanishing on the chosen scheme fibre, finite graded modules, flat base change and the theorem on formal functions.
\end{itemize}\medskip
\ifdefined\hypertarget\hypertarget{sourcepointer-1}{}\fi
\noindent\textit{Source reference.} Original equation 02NA: focused-seeds/AGP/sources/stacks-snapshot-20261011/tags.txt, source lines 3473--3473. Stable Stacks tag for an expressly allowed foundational or notational input; the counted construction remains the two concrete core obligations. The retained local core is not replaced by this pointer. See standard\_prerequisites.\par\medskip
\ifdefined\hypertarget\hypertarget{sourcepointer-2}{}\fi
\noindent\textit{Source reference.} Original lemma 073K: focused-seeds/AGP/sources/stacks-snapshot-20261011/tags.txt, source lines 9064--9064. Stable Stacks tag for an expressly allowed foundational or notational input; the counted construction remains the two concrete core obligations. The retained local core is not replaced by this pointer. See standard\_prerequisites.\par\medskip
\ifdefined\hypertarget\hypertarget{sourcepointer-3}{}\fi
\noindent\textit{Source reference.} Original theorem 08AU: focused-seeds/AGP/sources/stacks-snapshot-20261011/tags.txt, source lines 10544--10544. Stable Stacks tag for an expressly allowed foundational or notational input; the counted construction remains the two concrete core obligations. The retained local core is not replaced by this pointer. See standard\_prerequisites.\par\medskip
\ifdefined\hypertarget\hypertarget{sourcepointer-4}{}\fi
\noindent\textit{Source reference.} Original lemma 08AM: focused-seeds/AGP/sources/stacks-snapshot-20261011/tags.txt, source lines 10537--10537. Stable Stacks tag for an expressly allowed foundational or notational input; the counted construction remains the two concrete core obligations. The retained local core is not replaced by this pointer. See standard\_prerequisites.\par\medskip
\ifdefined\hypertarget\hypertarget{sourcepointer-5}{}\fi
\noindent\textit{Source reference.} Original lemma 0897: focused-seeds/AGP/sources/stacks-snapshot-20261011/tags.txt, source lines 10487--10487. Stable Stacks tag for an expressly allowed foundational or notational input; the counted construction remains the two concrete core obligations. The retained local core is not replaced by this pointer. See standard\_prerequisites.\par\medskip
\ifdefined\hypertarget\hypertarget{sourcepointer-6}{}\fi
\noindent\textit{Source reference.} Original lemma 0D2U: focused-seeds/AGP/sources/stacks-snapshot-20261011/tags.txt, source lines 16394--16394. Stable Stacks tag for an expressly allowed foundational or notational input; the counted construction remains the two concrete core obligations. The retained local core is not replaced by this pointer. See standard\_prerequisites.\par\medskip
\ifdefined\hypertarget\hypertarget{sourcepointer-7}{}\fi
\noindent\textit{Source reference.} Original theorem 08AZ: focused-seeds/AGP/sources/stacks-snapshot-20261011/tags.txt, source lines 10549--10549. Stable Stacks tag for an expressly allowed foundational or notational input; the counted construction remains the two concrete core obligations. The retained local core is not replaced by this pointer. See standard\_prerequisites.\par\medskip
\ifdefined\hypertarget\hypertarget{sourcepointer-8}{}\fi
\noindent\textit{Source reference.} Original lemma 08AS: focused-seeds/AGP/sources/stacks-snapshot-20261011/tags.txt, source lines 10542--10542. Stable Stacks tag for an expressly allowed foundational or notational input; the counted construction remains the two concrete core obligations. The retained local core is not replaced by this pointer. See standard\_prerequisites.\par\medskip
\ifdefined\hypertarget\hypertarget{sourcepointer-9}{}\fi
\noindent\textit{Source reference.} Original lemma 00DV: focused-seeds/AGP/sources/stacks-snapshot-20261011/tags.txt, source lines 512--512. Stable Stacks tag for an expressly allowed foundational or notational input; the counted construction remains the two concrete core obligations. The retained local core is not replaced by this pointer. See standard\_prerequisites.\par\medskip
\ifdefined\hypertarget\hypertarget{sourcepointer-10}{}\fi
\noindent\textit{Source reference.} Original lemma 08AY: focused-seeds/AGP/sources/stacks-snapshot-20261011/tags.txt, source lines 10548--10548. Stable Stacks tag for an expressly allowed foundational or notational input; the counted construction remains the two concrete core obligations. The retained local core is not replaced by this pointer. See standard\_prerequisites.\par\medskip
\begin{lemma}
\label{lemma-ample-on-fibre}
Let $Y$ be a Noetherian scheme. Let $X$ be an algebraic space over $Y$
such that the structure morphism $f : X \to Y$ is proper.
Let $\mathcal{L}$ be an invertible $\mathcal{O}_X$-module.
Let $\mathcal{F}$ be a coherent $\mathcal{O}_X$-module.
Let $y \in Y$ be a point such that $X_y$ is a scheme and
$\mathcal{L}_y$ is ample on $X_y$.
Then there exists a $d_0$ such that for all $d \geq d_0$ we have
$$
R^pf_*(\mathcal{F} \otimes_{\mathcal{O}_X} \mathcal{L}^{\otimes d})_y = 0
\text{ for }p > 0
$$
and the map
$$
f_*(\mathcal{F} \otimes_{\mathcal{O}_X} \mathcal{L}^{\otimes d})_y
\longrightarrow
H^0(X_y, \mathcal{F}_y \otimes_{\mathcal{O}_{X_y}} \mathcal{L}_y^{\otimes d})
$$
is surjective.
\end{lemma}

\begin{proof}
Note that $\mathcal{O}_{Y, y}$ is a Noetherian local ring.
Consider the canonical morphism
$c : \Spec(\mathcal{O}_{Y, y}) \to Y$, see
Schemes, Equation (\ref{schemes-equation-canonical-morphism}).
This is a flat morphism as it identifies local rings.
Denote momentarily $f' : X' \to \Spec(\mathcal{O}_{Y, y})$
the base change of $f$ to this local ring. We see that
$c^*R^pf_*\mathcal{F} = R^pf'_*\mathcal{F}'$ by
Cohomology of Spaces, Lemma
\ref{spaces-cohomology-lemma-flat-base-change-cohomology}.
Moreover, the fibres $X_y$ and $X'_y$ are identified.
Hence we may assume that $Y = \Spec(A)$ is the spectrum of
a Noetherian local ring $(A, \mathfrak m, \kappa)$ and $y \in Y$
corresponds to $\mathfrak m$. In this case
$R^pf_*(\mathcal{F} \otimes_{\mathcal{O}_X} \mathcal{L}^{\otimes d})_y =
H^p(X, \mathcal{F} \otimes_{\mathcal{O}_X} \mathcal{L}^{\otimes d})$
for all $p \geq 0$. Denote $f_y : X_y \to \Spec(\kappa)$ the projection.

\medskip\noindent
Let $B = \text{Gr}_\mathfrak m(A) =
\bigoplus_{n \geq 0} \mathfrak m^n/\mathfrak m^{n + 1}$.
Consider the sheaf $\mathcal{B} = f_y^*\widetilde{B}$
of quasi-coherent graded $\mathcal{O}_{X_y}$-algebras.
We will use notation as in Cohomology of Spaces, Section
\ref{spaces-cohomology-section-theorem-formal-functions}
with $I$ replaced by $\mathfrak m$.
Since $X_y$ is the closed subspace of $X$ cut out by
$\mathfrak m\mathcal{O}_X$ we may think of
$\mathfrak m^n\mathcal{F}/\mathfrak m^{n + 1}\mathcal{F}$
as a coherent $\mathcal{O}_{X_y}$-module, see
Cohomology of Spaces, Lemma \ref{spaces-cohomology-lemma-i-star-equivalence}.
Then
$\bigoplus_{n \geq 0} \mathfrak m^n\mathcal{F}/\mathfrak m^{n + 1}\mathcal{F}$
is a quasi-coherent graded $\mathcal{B}$-module of finite type
because it is generated in degree zero over $\mathcal{B}$
abd because the degree zero part is
$\mathcal{F}_y = \mathcal{F}/\mathfrak m \mathcal{F}$
which is a coherent $\mathcal{O}_{X_y}$-module.
Hence by Cohomology of Schemes, Lemma
\ref{coherent-lemma-graded-finiteness} part (2)
there exists a $d_0$ such that
$$
H^p(X_y, \mathfrak m^n \mathcal{F}/ \mathfrak m^{n + 1}\mathcal{F}
\otimes_{\mathcal{O}_{X_y}} \mathcal{L}_y^{\otimes d}) = 0
$$
for all $p > 0$, $d \geq d_0$, and $n \geq 0$. By
Cohomology of Spaces, Lemma
\ref{spaces-cohomology-lemma-relative-affine-cohomology}
this is the same as the statement that
$
H^p(X, \mathfrak m^n \mathcal{F}/ \mathfrak m^{n + 1}\mathcal{F}
\otimes_{\mathcal{O}_X} \mathcal{L}^{\otimes d}) = 0
$
for all $p > 0$, $d \geq d_0$, and $n \geq 0$.

\medskip\noindent
Consider the short exact sequences
$$
0 \to \mathfrak m^n\mathcal{F}/\mathfrak m^{n + 1} \mathcal{F}
\to \mathcal{F}/\mathfrak m^{n + 1} \mathcal{F}
\to \mathcal{F}/\mathfrak m^n \mathcal{F} \to 0
$$
of coherent $\mathcal{O}_X$-modules. Tensoring with $\mathcal{L}^{\otimes d}$
is an exact functor and we obtain short exact sequences
$$
0 \to
\mathfrak m^n\mathcal{F}/\mathfrak m^{n + 1} \mathcal{F}
\otimes_{\mathcal{O}_X} \mathcal{L}^{\otimes d}
\to \mathcal{F}/\mathfrak m^{n + 1} \mathcal{F}
\otimes_{\mathcal{O}_X} \mathcal{L}^{\otimes d}
\to \mathcal{F}/\mathfrak m^n \mathcal{F}
\otimes_{\mathcal{O}_X} \mathcal{L}^{\otimes d} \to 0
$$
Using the long exact cohomology sequence and the vanishing above
we conclude (using induction) that
\begin{enumerate}
\item $H^p(X, \mathcal{F}/\mathfrak m^n \mathcal{F}
\otimes_{\mathcal{O}_X} \mathcal{L}^{\otimes d}) = 0$
for all $p > 0$, $d \geq d_0$, and $n \geq 0$, and
\item $H^0(X, \mathcal{F}/\mathfrak m^n \mathcal{F}
\otimes_{\mathcal{O}_X} \mathcal{L}^{\otimes d}) \to
H^0(X_y, \mathcal{F}_y \otimes_{\mathcal{O}_{X_y}} \mathcal{L}_y^{\otimes d})$
is surjective for all $d \geq d_0$ and $n \geq 1$.
\end{enumerate}
By the theorem on formal functions
(Cohomology of Spaces, Theorem
\ref{spaces-cohomology-theorem-formal-functions})
we find that the $\mathfrak m$-adic completion of
$H^p(X, \mathcal{F} \otimes_{\mathcal{O}_X} \mathcal{L}^{\otimes d})$
is zero for all $d \geq d_0$ and $p > 0$.
Since $H^p(X, \mathcal{F} \otimes_{\mathcal{O}_X} \mathcal{L}^{\otimes d})$
is a finite $A$-module by
Cohomology of Spaces, Lemma
\ref{spaces-cohomology-lemma-proper-over-affine-cohomology-finite}
it follows from Nakayama's lemma (Algebra, Lemma \ref{algebra-lemma-NAK})
that $H^p(X, \mathcal{F} \otimes_{\mathcal{O}_X} \mathcal{L}^{\otimes d})$
is zero for all $d \geq d_0$ and $p > 0$.
For $p = 0$ we deduce from
Cohomology of Spaces, Lemma
\ref{spaces-cohomology-lemma-ML-cohomology-powers-ideal} part (3)
that $H^0(X, \mathcal{F} \otimes_{\mathcal{O}_X} \mathcal{L}^{\otimes d}) \to
H^0(X_y, \mathcal{F}_y \otimes_{\mathcal{O}_{X_y}} \mathcal{L}_y^{\otimes d})$
is surjective, which gives the final statement of the lemma.
\end{proof}

\bibliographystyle{plain}
\bibliography{my}
\end{document}
